\section{Methodology}
\label{sec:method}

\subsection{Questions and predictions}
We wrote the following predictions down before running any experiment
(\texttt{PLAN.md} in the artifact, dated 2026-11-23).
\begin{description}
\item[RQ1, where the time goes.] For long traces, Phase~1 dominates (we predicted
  10--30\,\textmu s per trace line); for many traced files, the per-file
  subprocesses of Phase~2 dominate (at least 3\,ms per file), and the \texttt{find}
  part grows with the number of files in the tree.
\item[RQ2, attribution.] Almost all of the gain over upstream comes from removing
  the per-file subprocesses and the per-file tree walk, not from threads; a
  single-threaded rewrite beats every parallel configuration on a fast device.
\item[RQ3, latency models.] With a warm cache (the normal case right after
  tracing), an extent lookup costs microseconds on all four file systems we test.
  With a cold cache, files created together share inode-table reads, so the mean
  number of device reads per file is far below one. A block-layer delay that lets
  requests overlap bounds thread speedup from above; a device that serves one
  request at a time gives no speedup; a per-file \texttt{sleep} overstates both.
\item[RQ4, equivalence.] Sequential and parallel modes disagree on unaligned
  requests, on file systems without FIEMAP, on unrelated \texttt{*.txt} files in the
  working directory, and above the parallel engine's queue size.
\end{description}

\subsection{The ablation ladder}
Every configuration reads the same \texttt{trace.result} in a fresh working
directory and scans the same file system mounted at \texttt{/mnt}:
\begin{description}
\item[(a)] upstream \texttt{processing.py}, unmodified (commit 796df86);
\item[(b)] the fork's sequential path (no flag; commit c162012);
\item[(c)] our single-threaded rewrite (\texttt{fastpath.py}, about 200 lines):
  Phase~1 in memory, one directory walk with \texttt{os.scandir} to build the
  inode map, each \texttt{<inode>.txt} written once, no subprocesses. It reproduces
  upstream's arithmetic, including its end-alignment formula, and the order GNU
  \texttt{sort -g} produces in the C locale; its output is byte-identical to (a) in
  every case we ran (\S\ref{sec:correct});
\item[(d)] the fork's \texttt{--parallel} path with 1, 4, 8 and 16 threads;
\item[(e)] (d) with 8 threads and the extent lookup replaced by a stub that returns
  an empty list, which isolates the cost of the lookup whose result is discarded.
\end{description}
Rung (c) plays the role of McSherry et al.'s competent single-threaded
baseline~\cite{mcsherry2015cost}: a parallel version is only interesting if it beats
it. For the whole pipeline we also compare upstream's three steps
(\texttt{processing.py}, \texttt{merge.py}, \texttt{hotness.sh}) with the rewrite
doing all three in one process.

\subsection{Workloads and traces}
We capture real traces with upstream's unmodified \texttt{trace.sh} on an ext4 file
system: upstream's own \texttt{read\_seq} and \texttt{read\_stride} benchmarks on a
1\,GiB file (128\,KiB \texttt{O\_DIRECT} reads), \texttt{tar} over a tree of about
5{,}000 source files (buffered reads of many files), a file-server-like
\texttt{fio} job (1{,}000 files of 1\,MiB, 4\,KiB random buffered reads), SQLite
point lookups on a 180\,MB database with a 1\,MiB page cache, and a single
\texttt{dd} read at a 5\,GiB offset. For controlled sweeps we generate traces in the
format \texttt{parse.sh} produces, for $F$ traced files chosen at random from a tree
of $T$ files.

\subsection{Slow devices}
We compare three ways of making storage slow, all with the same delay $d$:
\begin{description}
\item[sleep:] the fork's own method, a \texttt{time.sleep(d)} before every extent
  lookup, on a fast device with a warm cache;
\item[dm-delay:] the device-mapper delay target~\cite{kerneldmdelay} below the file
  system, so that only reads that reach the device are delayed and any number of
  them can wait concurrently; caches are dropped and the file system remounted
  before every cold run;
\item[QD1:] an NBD device served by nbdkit~\cite{nbdkit} with its delay filter and a
  single thread, so that requests are served one at a time, as by a disk with one
  actuator and no command queueing.
\end{description}
We count device reads from \texttt{/sys/block/<dev>/stat}. Reads are independent
of the host: they are what a real device would have to serve.

\subsection{Platform and statistics}
All runs use one laptop (AMD Ryzen~5 6600H, 6 cores and 12 threads, 15\,GB RAM, one
NVMe SSD) running Linux 6.18 in WSL2 (Ubuntu 24.04, Python 3.12.3, GNU findutils and
coreutils from Ubuntu). File systems under test live on loop devices with direct
I/O enabled, backed by a file on the WSL2 virtual disk. Unmodified upstream code
hard-codes \texttt{/mnt}, so every run executes in a private mount namespace in
which the file system under test is mounted at \texttt{/mnt}. Each configuration
runs ten times unless stated otherwise, with the order of configurations shuffled
in every repetition; we report medians with 95\% percentile-bootstrap intervals
(10{,}000 resamples) and compute intervals for speedups by resampling numerator and
denominator independently~\cite{hoefler2015scientific}. Raw results are JSON lines
with the commit hashes, kernel and configuration of every run; every table and
figure is generated from them by one script.
